\chapter{C14.4 at large $\zeta$: every unmatched Bob residue}
\label{component:p1a6cc124f5db}
\noindent\textbf{Component }\path{p1a6cc124f5db}\par
{\small\noindent Source: \path{BASE/inherited_r19/inherited_r18/proofs/c1_large_eta/off_matched_bob.tex}\par}
{\small\noindent SHA256: \path{1a6cc124f5db69e647dbc27b8df98ba0f8cbbf3f311a8504998f074d4fa265f0}\par}
{\small\noindent Inventory use: EB1--EB10 | EB exact homogeneous parent behind SE | EB1-EB10 complete homogeneous full-state unmatched chart\par}
\medskip
\section{C14.4 at large $\zeta$: every unmatched Bob residue}
\label{p1a6cc124f5db:sec:r18-c144-bob-residue}

\paragraph{Precise admission.}
Retain the entire original C14.4 source: $x,L>0$ infinitesimal,
$\eta$ unbounded, $\zeta=L/x$ noninfinitesimal, all original physical
intersections, every positive opening, and either permitted literal
$D=0$ or $\beta=0$. This section additionally assumes $\zeta$
unbounded. Put
\[
 \sigma=1/\eta,\qquad S=\sigma+L,\qquad
 \nu=\sigma/S,\qquad \ell=L/S=1-\nu.
\]
For a real projective Bob test choose its normalized homogeneous
coordinates $(c,s)$, so $c^2+s^2=1$; on the finite chart these are
$c=(1+v^2)^{-1/2}$, $s=vc$. The complete original parent is tensor SOS
at every independent test and state satisfying
\[
                         \overline{s^2}\ne\overline\ell.       \tag{EB1}
\]
Bars mean real residues of bounded elements in their respective fields.
This compares real residues only, not independent infinitesimal rates.
The adjacent-$C,A$ theorem already covers noninfinitesimal Alice
tests. Hence only $u$ infinitesimal is proved below. No all-tests
closure of C14.4 is asserted: the matched residue and bounded-unit
$\zeta$ source remain.

\paragraph{Source limits, including both sphere endpoints.}
Use ES1--ES10 without changing their source entry. Since
$x/S\le x/L=1/\zeta$ is infinitesimal, ES7 gives
$e=xh/\sigma$ infinitesimal, and hence $E\sim1$.
Moreover $M=1+xr-\theta/\zeta\sim1$. ES8 then gives $m$ infinitesimal.
The exact equation ES6, divided by $S$, is
\[
 e(B/S)=(x/S)H_0+\nu\mathcal H b_0^2.
\]
Both summands are nonnegative; $B/S$ and $\mathcal H$ tend to one.
Consequently $\nu b_0^2$ is infinitesimal. Therefore ES5 yields
\[
 \frac DS-\frac{\nu^2}{2}\in\operatorname{Inf},\qquad
 \frac\alpha S-\frac{\ell^2}{2}\in\operatorname{Inf},
 \qquad m\in\operatorname{Inf}_{>0}.                         \tag{EB2}
\]
For the second formula, divide the exact ES9 expression by $S$:
\[
 \frac\alpha S=
 \frac{(1+\sigma r)\{\ell R e+\ell^2ME\}}
 {2Q(P/S)}.
\]
All displayed denominator factors tend to one. These are additive
residue assertions, valid also when $\nu$ or $\ell$ is infinitesimal.
In particular no division by either sphere coordinate is used.

The opening-cleared ES8 bounds give
\[
 c_*=SC(1-r)\sim\tfrac12,\qquad
 a_*=SAr\sim\tfrac12.
\]
Define the pure-source harmonic quantities
\[
 H=S(C+A),\qquad \rho=C/(C+A),\qquad
 a_h=SCA/(C+A)=\frac{c_*a_*}{c_*r+a_*(1-r)}.
                                                               \tag{EB3}
\]
Then $1/\sqrt H$ is bounded, $0<\rho<1$, and $a_h\sim1/2$,
uniformly through both $r$ openings. No upper bound on $C$ or $A$
has been assumed.

\paragraph{An exact all-state frame for every Bob chart.}
Let $D_u=1+u^2$ and $A_u=1+u^2s^2$, both positive testing units.
Use the separated roots
\[
 p=\frac{uW}{\sqrt S},\quad q=\sqrt S Z,\quad
 y=\frac Y{\sqrt S},\quad N_y=\sqrt H(y-\rho p),
 \quad N_m=\sqrt m\,sW,
\]
\[
 N_b=\sqrt{A_u}\left(X+\frac{uc^2}{A_u}Z\right),
 \qquad N_r=\frac{\sqrt{D_u}\,usc}{\sqrt{A_u}}Z.
                                                               \tag{EB4}
\]
Only pure-source and pure-testing radicals or inverses occur.
In particular no inverse of $c$, $s$, $u$, or a state coordinate is
needed. The inverse harmonic allocation
$y=\rho p+N_y/\sqrt H$ has bounded coefficients.

Write $b=X+uZ$. The exact original signed core $K/D_v$ in finite
Bob coordinates equals the following homogeneous expression:
\[
\begin{split}
 K_B={}&N_b^2+N_r^2-2D_usc\,pq
       +Sc^2(p-y)^2+SD_us^2y^2\\
 &+\frac{2\sqrt S\,usc}{\sqrt{A_u}}(p-y)N_b
  +\frac{2\sqrt S\,us^2\sqrt{D_u}}{\sqrt{A_u}}(p-y)N_r.
\end{split}                                                    \tag{EB5}
\]
Here $D_us$ denotes $D_u\,s$, and similarly in the other products.
One derivation starts with CE2 divided by $D_v$, sets
$R=usZ$, and completes its entire $(b,R)$ block:
\[
 A_ub^2-2D_us\,bR+D_uR^2=N_b^2+N_r^2.
\]
The cross term from
$[\sqrt S\,c(p-y)+usb]^2$ is allocated by
$cb=cN_b/\sqrt{A_u}+\sqrt{D_u}sN_r/\sqrt{A_u}$.
This proves EB5 without deleting a first state normal.

The original positive event roots have the exact bounded expressions
\[
 \begin{split}
 M_B={}&N_m-\frac{\sqrt m\,c}{\sqrt{A_u}}N_b
       -\frac{\sqrt{mD_u}\,s}{\sqrt{A_u}}N_r
       +\sqrt{mS}\,us\,y,\\
 D_B={}&\sqrt{DS}\,c(p-y)
       +\frac{\sqrt D\,us}{\sqrt{A_u}}N_b
       -\sqrt{D/S}\,\frac{D_us}{A_u}q,\\
 A_B={}&\sqrt{\alpha D_uS}\,s y
       -\sqrt{\alpha D_u/S}\,c q.
 \end{split}                                                   \tag{EB6}
\]
They are, respectively, the normalized original $m,D,\alpha$ event
amplitudes. The formulas are valid at $D=0$ as well. In the $D$ row,
the identity
$usb-D_usZ=usN_b/\sqrt{A_u}-D_usZ/A_u$
avoids an apparent reciprocal-$c$ pole.

Consequently the full original quarter-parent, divided by $D_v$ on
the finite chart, is exactly
\[
 \mathcal P_B=N_y^2+a_hp^2+ASu^2y^2
              +M_B^2+D_B^2+A_B^2+\beta b^2+\tfrac12K_B.
                                                               \tag{EB7}
\]
The two explicit positive rows $ASu^2y^2$ and $\beta b^2$ can be
retained unchanged. All other rows in EB7 use EB4 once. The formula
is a full four-state identity and is not a restriction to a state ray.

\paragraph{Strictness precisely off the matched residue.}
Every coefficient of EB5 except the printed $pq$ term is either
its fixed positive pivot or infinitesimal after EB3. In EB6, the
$m$ row is $N_m$ plus infinitesimal coefficients; the $D$ and
$\alpha$ rows have only their $q$ coefficients at residue.
Use EB2, $A_u,D_u\sim1$, and EB3. The resulting residue Gram on
$(N_y,N_m,N_b,N_r,p,q)$ is the direct sum of
$\operatorname{diag}(1,1,1/2,1/2)$ and
\[
 \frac12
 \begin{pmatrix}
 1&-\bar s\bar c\\
 -\bar s\bar c&\bar\nu^2\bar s^2+\bar\ell^2\bar c^2
 \end{pmatrix}.                                                \tag{EB8}
\]
The determinant of this last block is
\[
             \frac14(\bar s^2-\bar\ell)^2.                    \tag{EB9}
\]
Thus EB1 makes the entire bounded finite Gram strictly positive.
The separated bounded-Gram lemma gives tensor SOS, and restoring the
two retained positive rows proves the complete parent. The roots are
allowed to be dependent or even identically zero at an endpoint; the
claim concerns the coefficient Gram, so this does not affect the
argument.

The same homogeneous formulas allow $c=0$ and $s=0$ literally.
They include every Bob projective endpoint, every sign, and every
local Alice axis. Nonlocal Alice charts use the inherited adjacent
$C,A$ theorem. The real and imaginary state parts are treated
separately. Every normalization is pure source or pure testing.

\paragraph{The exact remaining obstacle.}
The leading scalar has the useful two-square identity
\[
 \tfrac12p^2-sc\,pq
 +\tfrac12(\nu^2s^2+\ell^2c^2)q^2
 =\tfrac12(p-scq)^2+\tfrac12(s^2-\ell)^2q^2.                  \tag{EB10}
\]
It does not complete the matched-residue chart: there the second
pivot collapses, and the actual source corrections must be allocated
through exact contact conductors. Arbitrary small positive $D$ or
$\beta$ openings do not supply a missing fixed pivot. Unit $\zeta$
also requires new source limits because $e$ can have nonzero residue.
These two obligations are retained, not declared empty.

The focused checker verifies EB5--EB7 as complete homogeneous
quadratic identities and EB9--EB10 algebraically. Source estimates
EB2--EB3 are the analytical consequences of the already proved
ES identities given above. This is a partial-domain admission only.
